\documentclass[11pt]{article}

\usepackage[a4paper,margin=1in]{geometry}
\usepackage{amsmath,amssymb,mathtools}
\usepackage[hidelinks]{hyperref}

\title{Global Foliation Trace Route}
\author{}
\date{}

\begin{document}
\maketitle

This note records the sharp version of Plan 2: do not select one good
near-boundary level. Use the whole outer family of level sets. The point is
that the Holder/Serrin loss \(D_H\) is not just extra pointwise information; in
level-set coordinates it is the radial kinetic energy of the foliation.
Together with the isoperimetric loss \(D_I\), it gives a trace estimate for the
boundary asymmetry.

The argument below is written first in the coherent near-spherical foliation
regime. This is the part that should replace the final selected-minimizer
compactness step. The remaining task is an entry theorem saying that a small
global level-set deficit puts the actual torsion levels into this weak
foliation gauge on a fixed outer annulus.

Throughout this note normalize
\[
|\Omega|=\omega_n,\qquad B=B_1,
\]
and write
\[
\delta_T(\Omega)=E(\Omega)-E(B).
\]
Let \(u=u_\Omega\) be the torsion function,
\[
E_t=\{u>t\},\qquad m(t)=|E_t|,
\]
and introduce the volume radius
\[
\rho=\left(\frac{m(t)}{\omega_n}\right)^{1/n}.
\]
Equivalently, write \(t=t(\rho)\) and \(E_\rho=E_{t(\rho)}\). The exact
level-set identity is
\[
2\delta_T(\Omega)
=
\int_0^{\|u\|_\infty}
\frac{D_H(t)+D_I(t)}
{n^2\omega_n^{2/n}m(t)^{1-2/n}}\,dt.
\]

\section{Coherent Outer Foliation Gauge}

Fix \(0<\rho_*<1\), for instance \(\rho_*=1/2\). Assume that for a.e.
\(\rho\in[\rho_*,1]\) the regular level surface \(\partial E_\rho\) can be
written in the modulated graph gauge
\[
\partial E_\rho
=
\left\{
z(\rho)+(\rho+h(\rho,\theta))\theta:
\theta\in\partial B
\right\},
\]
with
\[
\int_{\partial B} h(\rho,\theta)\,d\theta=0,
\qquad
\int_{\partial B} h(\rho,\theta)\theta_i\,d\theta=0
\quad(i=1,\ldots,n).
\]
The zero mode is removed by using the volume radius \(\rho\). The first modes
are removed by the moving center \(z(\rho)\). This is important: translation is
a neutral mode for Fraenkel asymmetry, so the center is allowed to move with
\(\rho\).

Define the non-neutral quadratic form
\[
Q(h)
:=
\int_{\partial B}
\left(|\nabla_\theta h|^2-(n-1)h^2\right)\,d\theta.
\]
On the subspace orthogonal to constants and first spherical harmonics,
\[
Q(h)\ge c_n\|h\|_{H^1(\partial B)}^2.
\]

The coherent-foliation estimate needed from the geometry is
\[
\boxed{
\int_{\rho_*}^1
\left(
\rho^{n+1}\|\partial_\rho h(\rho)\|_{L^2(\partial B)}^2
+\rho^{n-1}Q(h(\rho))
\right)\,d\rho
\le
C_n\delta_T(\Omega).
}
\tag{F}
\]
Once (F) is known, sharp stability follows immediately.

\section{Trace Estimate: No Boundary-Slab Loss}

Since \(\rho\in[\rho_*,1]\), the weights \(\rho^{n+1}\) and \(\rho^{n-1}\) are
bounded above and below by dimensional constants. For every non-neutral mode,
the one-dimensional trace theorem gives
\[
\|h(1)\|_{L^2(\partial B)}^2
\le
C_{\rho_*}
\int_{\rho_*}^1
\left(
\|\partial_\rho h(\rho)\|_{L^2(\partial B)}^2
+\|h(\rho)\|_{L^2(\partial B)}^2
\right)\,d\rho.
\]
Using the coercivity of \(Q\) on the chosen gauge,
\[
\boxed{
\|h(1)\|_{L^2(\partial B)}^2
\le
C_n
\int_{\rho_*}^1
\left(
\rho^{n+1}\|\partial_\rho h(\rho)\|_2^2
+\rho^{n-1}Q(h(\rho))
\right)\,d\rho.
}
\tag{T}
\]
Combining (T) with (F) gives
\[
\|h(1)\|_{L^2(\partial B)}^2\le C_n\delta_T(\Omega).
\]
For a near-spherical set in the same volume and translation gauge,
\[
\mathcal A(\Omega)^2\le C_n\|h(1)\|_{L^2(\partial B)}^2.
\]
Hence
\[
\boxed{
\mathcal A(\Omega)^2\le C_n\delta_T(\Omega).
}
\]

This is the sharp Saint--Venant exponent. No parameter \(\eta\) is introduced:
the trace is taken from a fixed annulus, and the radial derivative term
supplied by \(D_H\) prevents the boundary deformation from being hidden in an
arbitrarily thin layer.

\section{Why \(D_I\) Gives the Angular Part}

In the graph gauge,
\[
P(E_\rho)^2-P(B_\rho)^2
\ge
c_n\rho^{2n-4}Q(h(\rho))
-\hbox{higher order}.
\]
This is the Fuglede expansion of the perimeter deficit on the level
\(\partial E_\rho\), after removing the volume and translation modes.

Since \(m(t(\rho))=\omega_n\rho^n\), the denominator in the exact identity
satisfies
\[
n^2\omega_n^{2/n}m(t(\rho))^{1-2/n}\simeq_n \rho^{n-2}.
\]
In a near-spherical foliation,
\[
-t_\rho(\rho)\simeq_n \rho.
\]
Therefore
\[
\int_{\rho_*}^1
\frac{D_I(t(\rho))}
{n^2\omega_n^{2/n}m(t(\rho))^{1-2/n}}
(-t_\rho(\rho))\,d\rho
\ge
c_n
\int_{\rho_*}^1
\rho^{n-1}Q(h(\rho))\,d\rho
-\hbox{errors}.
\]
For a sufficiently small graph radius, the errors are absorbed by the left-hand
side or by the \(D_H\) contribution.

\section{Why \(D_H\) Gives the Radial Kinetic Part}

Let
\[
X(\rho,\theta)=z(\rho)+(\rho+h(\rho,\theta))\theta
\]
parametrize \(\partial E_\rho\), and let \(\nu_\rho\) be the outer unit normal
to \(E_\rho\). Define the normal velocity of the level foliation with respect
to the volume radius:
\[
V(\rho,\theta):=\partial_\rho X(\rho,\theta)\cdot\nu_\rho(\theta).
\]
Since \(u(X(\rho,\theta))=t(\rho)\) and the outward normal is
\(\nu_\rho=-\nabla u/|\nabla u|\),
\[
|\nabla u|(X(\rho,\theta))
=
\frac{-t_\rho(\rho)}{V(\rho,\theta)}.
\]
Thus the scalar factor \(-t_\rho\) cancels out of the Holder deficit:
\[
D_H(t(\rho))
=
\left(\int_{\partial E_\rho}V\,d\mathcal H^{n-1}\right)
\left(\int_{\partial E_\rho}\frac1V\,d\mathcal H^{n-1}\right)
-P(E_\rho)^2.
\]
Moreover, because \(\rho\) is the volume radius,
\[
\int_{\partial E_\rho}V\,d\mathcal H^{n-1}
=
\frac{d}{d\rho}|E_\rho|
=
n\omega_n\rho^{n-1}
=P(B_\rho).
\]
So \(D_H\) is exactly the Cauchy deficit for the normal velocity \(V\). In the
graph gauge,
\[
V=1+z'(\rho)\cdot\theta+\partial_\rho h+\hbox{quadratic terms}.
\]
The moving center absorbs the first spherical harmonics, and the Cauchy deficit
controls the variance of \(V\). Projecting to the non-neutral modes gives
\[
D_H(t(\rho))
\ge
c_n\rho^{2n-2}
\|\partial_\rho h(\rho)\|_{L^2(\partial B)}^2
-\hbox{lower order errors}.
\]
After multiplying by the exact identity weight,
\[
\int_{\rho_*}^1
\frac{D_H(t(\rho))}
{n^2\omega_n^{2/n}m(t(\rho))^{1-2/n}}
(-t_\rho(\rho))\,d\rho
\ge
c_n
\int_{\rho_*}^1
\rho^{n+1}\|\partial_\rho h(\rho)\|_2^2\,d\rho
-\hbox{errors}.
\]
This is the term missing from the one-level replacement argument.

\section{The Sharp Conditional Theorem}

\textbf{Theorem.}
Fix \(\rho_*\in(0,1)\). There are constants
\(\varepsilon_0,c,C>0\), depending only on \(n\) and \(\rho_*\), with the
following property. Suppose the outer level sets of \(\Omega\) admit the
coherent foliation gauge above on \([\rho_*,1]\), with graph size at most
\(\varepsilon_0\), and suppose the quadratic lower bounds in Sections 3 and 4
hold with absorbable errors. Then
\[
\mathcal A(\Omega)^2\le C\delta_T(\Omega).
\]

\emph{Proof.}
Add the lower bounds from Sections 3 and 4 and use the exact level-set identity
restricted to the interval \(\rho\in[\rho_*,1]\). This gives (F). The trace
estimate (T) gives
\[
\|h(1)\|_2^2\le C\delta_T(\Omega).
\]
Finally, the standard near-spherical comparison between the Fraenkel asymmetry
and the \(L^2\) boundary graph, after volume and translation normalization,
gives
\[
\mathcal A(\Omega)^2\le C\|h(1)\|_2^2.
\]

\section{What Must Replace the Selection Principle}

The selection principle should be replaced by the following level-set entry
statement, not by a stronger one-level Serrin theorem.

\textbf{Outer foliation entry lemma.}
If \(\delta_T(\Omega)\) is sufficiently small, then the actual torsion levels
\(E_\rho\), \(\rho\in[\rho_*,1]\), admit a modulated weak graph representation
\[
\partial E_\rho
=
z(\rho)+(\rho+h(\rho,\theta))\theta
\]
with
\[
h\in H^1_\rho L^2_\theta\cap L^2_\rho H^1_\theta
\]
on the non-neutral modes, and the action estimate (F) holds.

This is a quantitative theorem about the original level sets. It does not
construct a penalized minimizer. The proposed proof is:
\begin{enumerate}
\item Use the exact identity to control the total \(D_I+D_H\) action on the
fixed outer annulus \(\rho\in[\rho_*,1]\).
\item Use quantitative isoperimetry level by level to modulate each \(E_\rho\)
by a center \(z(\rho)\) and remove the zero and first modes.
\item Use the exact velocity formula for \(D_H\) to control the \(L^2_\rho\)
derivative of the non-neutral part of this modulation.
\item Apply the Hilbert-space trace theorem to recover the boundary graph.
\end{enumerate}

This is the precise sense in which the global level-set data can act as the
surrogate for the selection principle. The one-level route loses \(1/\eta\);
the foliation route replaces that average by radial kinetic energy.

\section{Immediate Next Proof Obligation}

The next concrete task is to prove the two quadratic lower bounds with all
remainders:
\[
D_I(t(\rho))
\gtrsim
\rho^{2n-4}Q(h(\rho)),
\]
and
\[
D_H(t(\rho))
\gtrsim
\rho^{2n-2}
\|\partial_\rho h_{\ge2}(\rho)\|_2^2
-C\rho^{2n-4}Q(h(\rho))\,\|h\|_{C^1}.
\]
These estimates are local computations in the moving graph gauge. Once written
carefully, the only genuinely global issue is the outer foliation entry lemma.

\end{document}
